% Study sheet: Volume II, Chapter 8, Fleet Orchestration.
% Numbers: \val{...} from fleet_orchestration_studysheet_numbers.tex, written by fleet_orchestration_studysheet_numbers.py.
\documentclass{studysheet}
\input{fleet_orchestration_studysheet_numbers}
\sheetvolume{2}
\sheetchapter{8}{Fleet Orchestration}
\sheetsources{Numbers from mlsysim and the chapter's scenarios; prices are scenario rates per GPU-hour; assumed inputs marked.}

\begin{document}
% ===================================================================== FRONT: LEARN
\frontside
\begin{bigidea}
A shared fleet wastes money when jobs wait beside idle accelerators. The \hl{scheduler} decides which
jobs run, where, and when one yields, and \hl{feasibility} comes before fairness. In C\textsuperscript{3}
terms, orchestration spends coordination to turn compute and communication into useful progress.
\end{bigidea}

\begin{leftcol}
\sheetsec{Key ideas}
\begin{ideas}
\idea{Feasibility before fairness.}{A share a job cannot run on delivers nothing: check gang, fit, topology
  and memory first, then share fairly.}
\idea{A rigid gang is all or nothing.}{Half a gang runs no steps; partial holds can deadlock. Grant
  atomically, reserve for the queue head, backfill jobs that end in time.}
\idea{Admit, place, reconcile.}{The scheduler's view is stale: reuse a lost node's GPUs only
  after its lease lapses and old work is fenced.}
\idea{Placement sets the communication term.}{Give the chattiest group the fastest tier; judge by
  completion time, including the wait for a compact slot.}
\idea{Every transition must repay itself.}{A preemption, resize or spot reclaim pays lost work, restart,
  reload and warmup; move only if the gain outlasts them.}
\idea{Count utilization in tiers.}{A full queue beside idle GPUs is usually a policy written for an
  older fleet, not a bug.}
\end{ideas}

\vspace{10pt}
\sheetsec{Where each group goes}
\begin{sheettable}{@{}l l L@{}}
\thead{Tier} & \thead{Per direction} & \thead{Put there}\\ \midrule
Node (NVLink) & \val{nvlink} GB/s & tensor parallel: every layer\\
One leaf switch & \val{ndr} GB/s & pipeline stage boundaries\\
Spine (oversubscribed) & \val{spinelo}--\val{spinehi} GB/s & data parallel: once per step\\
\end{sheettable}

\vspace{10pt}
\sheetsec{Which utilization tier is low?}
\begin{sheettable}{@{}l L@{}}
\thead{Tier counts} & \thead{If low, check}\\ \midrule
Allocated: whole gangs & demand, quotas, health; requested vs. needed\\
Active: hardware busy & input supply, communication waits\\
Productive: progress kept & replay, synchronization, overhead\\
\end{sheettable}

\vspace{10pt}
\sheetsec{Match the signal to the objective}
\begin{sheettable}{@{}l L l@{}}
\thead{Signal} & \thead{Answers} & \thead{Example}\\ \midrule
GPU-hours used & unknown run lengths & Tiresias\\
Iteration boundaries & costly preemption & Gandiva\\
Normalized finish time & unequal delay & Themis\\
Goodput & batch size with GPUs & Pollux\\
\end{sheettable}
\end{leftcol}%
\begin{rightcol}
\sheetsec{Formulas you will use}
\fcard{Queue wait (one server)}%
  {\frac{\QueueWaitTime}{1/\ServiceRate} = \under{\frac{\ServingUtilization}{1-\ServingUtilization}}{load}
   \cdot \under{\frac{1+\ServiceTimeVariationCoefficient^2}{2}}{variability}}%
  {Mean wait $\QueueWaitTime$ in mean-job durations $1/\ServiceRate$, at offered load $\ServingUtilization$ and
   job-length spread $\ServiceTimeVariationCoefficient$ (standard deviation over mean). Slowdown
   $\JobSlowdown_{\JobIndex} = (\JobQueueWaitTime_{\JobIndex}+\IsolatedJobRunTime_{\JobIndex})/\IsolatedJobRunTime_{\JobIndex}$,
   with isolated run $\IsolatedJobRunTime_{\JobIndex}$.}%
  {Poisson arrivals, one first-come server; ignores gangs and packing.}%
  {\begin{tikzpicture}[sheet]
     \draw[dim] (20,0) -- (230,0); \draw[dim] (20,0) -- (20,52);
     \foreach \x/\t in {0/0,0.5/{0.5},0.8/{0.8},1/1}{
       \draw[dim] ({20+\x*200},0) -- ({20+\x*200},-2.5);
       \node[small, anchor=north] at ({20+\x*200},-4) {\t};}
     \draw[accent, dashed, line width=0.5pt] (220,0) -- (220,52);
     \draw[muted, line width=0.9pt] plot[domain=0:0.98, samples=60, smooth] ({20+\x*200},{\x/(1-\x)});
     \draw[accent, line width=1.1pt] plot[domain=0:0.909, samples=60, smooth] ({20+\x*200},{\val{varmixed}*\x/(1-\x)});
     \filldraw[fill=accent, draw=accent] (180,\val{waitmixed}) circle (2pt);
     \filldraw[fill=white, draw=muted] (180,\val{waituniform}) circle (2pt);
     \node[lab, anchor=east] at (176,{\val{waitmixed}+5}) {mixed, $\ServiceTimeVariationCoefficient = \val{cvmixed}$: \val{waitmixed}$\times$};
     \node[note, anchor=west] at (24,46) {wait $\div$ mean job};
     \node[note, anchor=west] at (24,36) {gray: $\ServiceTimeVariationCoefficient = \val{cvuniform}$, \val{waituniform}$\times$ at \val{loadref}};
     \node[note, anchor=north] at (120,-12) {offered load};
   \end{tikzpicture}}
\fcard{Dominant share (DRF)}%
  {\DominantShare_{\TenantIndex} = \max_{\ResourceTypeIndex}\frac{\TenantResourceAllocation}{\FleetResourceCapacity}}%
  {A tenant's largest share of any resource: GPUs, CPU, memory, network. Progressive filling gives the
   next bundle to the tenant with the smaller $\DominantShare$ until a resource runs out; it equalizes
   $\DominantShare$, not every share.}%
  {Only over bundles that pass the feasibility tests.}{}
\fcard{Effective spot price}%
  {\EffectiveSpotTrainingCost = \SpotHourlyPrice \times \frac{\BilledAllocationTime}{\RetainedProductiveTime}}%
  {Price per hour of retained progress: billed time is progress kept plus checkpoints, restarts
   and discarded work. Spot pays while it stays below the on-demand price.}%
  {Scenario rates: \$\val{spotrate} spot, \$\val{ondemand} on demand, per GPU-hour.}{}
\end{rightcol}

% ===================================================================== BACK: PRACTICE
\backside
\begin{answers}
\ans{Finish}{\val{ybilled} $\div$ \val{yint} = \val{yratio}; \$\val{spotrate} $\times$ \val{yratio} = \$\val{yeff},
  \val{ysavepct}\% below on demand. Break-even: \$\val{spotrate}$(t+\val{ylosth})/t$ = \$\val{ondemand}
  gives $t$ = \val{ylosth} $\times$ \val{spotrate} $\div$ \val{gap} $\approx$ \val{beh} h, or \val{bemin}
  minutes (spot card).}
\ans{1}{(a) Load factor \val{loadfactor} becomes \val{qloadfactor}; times variability \val{varmixed}: \val{waitmixed}
  becomes \val{qwait} mean jobs. (b) $\IsolatedJobRunTime_{\JobIndex}$ is one mean job: \val{qwait} + one = \val{qslow} (queue card). (c) GPU-hours used: favor jobs with less (signals table).}
\ans{2}{(a) A: \val{ashare} of every resource; P: cores, \val{ccores}/\val{capcores} = \val{cshare}. (b) Two A
  bundles per P: A \val{abundles}, P \val{cbundles}, both $\DominantShare$ = \val{qdrfshare}. Cores fill
  (\val{acorestotal} + \val{ccorestotal}), so \val{qgpusidle} GPUs sit idle (DRF card).}
\columnbreak
\ans{3}{Allocated, in the older pools. Requested vs. measured peak memory per pool
  shows templates over-asking for the newest pool. Fix the policy (capability matching,
  backfill; tier table).}
\ans{4}{\val{tpgb} GB $\div$ \val{nvlink} GB/s = \val{tpnv} ms on NVLink; $\div$ \val{ndr} GB/s = \val{tpib} ms
  across nodes, \val{tpslow}$\times$ slower, paid every layer, every step. Keep tensor parallel
  in a node (placement table).}
\end{answers}

\begin{leftcol}
\sheetsec{Worked example}
\begin{practicebox}\boxtitle{Is spot cheaper for a \val{wdays}-day run on \val{wgpus} GPUs?}
\val{wmodel} needs \val{whours} h of training: \$\val{ondemand} per GPU-hour on demand, or
\$\val{spotrate} on spot with \val{ckptpct}\% checkpoint overhead and \val{perday} interruptions per training day
(\val{wint} in all), each \val{restarth} h to restart.\par\vspace{5pt}
{\centering
\begin{tikzpicture}[sheet]
  \pgfmathsetmacro\sc{140/\val{drawod}}
  \node[note, anchor=west] at (0,26) {on demand};
  \draw[nbox] (50,20) rectangle ({50+\val{drawod}*\sc},32);
  \node[lab, anchor=west] at ({54+\val{drawod}*\sc},26) {\$\val{wodk}k};
  \pgfmathsetmacro\a{50+\val{drawbase}*\sc} \pgfmathsetmacro\b{\a+\val{drawover}*\sc}
  \node[note, anchor=west] at (0,6) {spot};
  \draw[accbox] (50,0) rectangle (\a,12);
  \draw[hotbox] (\a,0) rectangle (\b,12);
  \node[lab, anchor=west] at ({\b+4},6) {\$\val{wspotk}k};
  \node[note, anchor=north west] at (50,-3) {solid: checkpoints and restarts};
\end{tikzpicture}\par}\vspace{6pt}
\begin{steps}
\item On demand: \val{wgpus} $\times$ \val{whours} h $\times$ \$\val{ondemand} = \res{\$\val{wodcost}}.
\item Billed spot hours: \val{whours} + \val{ckpth} (\val{ckptpct}\% checkpoints) + \val{wint} $\times$ \val{restarth}
  (restarts) = \val{wbilled} h.
\item Effective price: \$\val{spotrate} $\times$ \val{wbilled} $\div$ \val{whours} = \$\val{spotrate} $\times$
  \val{wratio} = \res{\$\val{weff}} per GPU-hour of progress.
\item Bill \val{wgpus} $\times$ \val{wbilled} $\times$ \$\val{spotrate} = \$\val{wspotcost}: \res{\val{wsavepct}\% less},
  but \val{wlongerpct}\% more GPU-hours.
\end{steps}
\takeaway{The discount survives while overhead stays small.}
\end{practicebox}

\vspace{10pt}
\sheetsec{Your turn: finish it}
\begin{practicebox}\boxtitle{Slow storage: does the discount survive?}
Each interruption bills a \val{yckpt}-minute save and a \val{yrestart}-minute restart, once per
\val{yint} h of progress kept.\par
\textbf{Given:} each event bills \val{yint} + \val{ylosth} = \val{ybilled} h for \val{yint} h kept.\par
\textbf{You:} find the effective price and saving against \$\val{ondemand}, then the progress per
interruption at which spot stops paying.\par\vspace{34pt}
\end{practicebox}
\end{leftcol}%
\begin{rightcol}
\sheetsec{Common mistakes}
\mistake{``Half a gang still makes half the progress.''}%
  {A rigid \val{gang}-GPU job holding \val{gangheld} runs no steps; an elastic one could start at its minimum gang. Two
   whole-cluster jobs each holding half deadlock: \val{cluster} GPUs $\times$ \$\val{ondemand} =
   \$\val{clusterhour} an hour, idle.}{}
\mistake{``\val{assigned}\% of GPUs are assigned, so \val{assigned}\% are working.''}%
  {\val{incomplete}\% sit in incomplete gangs and \val{inputwait}\% wait for input: \val{complete}\% is in
   complete groups, at most \val{executing}\% training.}{}
\mistake{``Size each team's quota to its peak.''}%
  {Peaks of \val{peaka} and \val{peakb} need \val{peaksum} GPUs; \val{shared} shared cover either
   peak (not both) if borrowed jobs yield in time. The extra \val{extra} (\val{extrapct}\%) idle
   whenever peaks miss.}{}

\vspace{4pt}
\sheetsec{Check yourself}
\begin{qs}
\item \assumed{} a mixed queue ($\ServiceTimeVariationCoefficient = \val{cvmixed}$) rises from \val{loadref} to
  \val{qload} offered load. (a) Mean wait, in mean jobs? (b) Slowdown of a mean-length job? (c) Which signal can cut short jobs' wait without runtime estimates?
\item \assumed{} the chapter's cluster: \val{capgpus} GPUs, \val{capcores} cores, \val{capmem} GB. A's replica:
  \val{agpu} GPU, \val{acores} cores, \val{amem} GB; P's preprocessing job: \val{cgpu} GPU, \val{ccores} cores, \val{cmem} GB. (a) Each bundle's dominant share?
  (b) Fill progressively: bundles each, and how many GPUs idle, why?
\item Pools of A100-80 GB, A100-40 GB and V100 show middling utilization; hundreds of jobs
  wait, all for A100-80 GB. Which tier is low, where? What separates a request problem from a demand
  problem?
\item A tensor-parallel group is split across two nodes under one leaf. How much slower is a \val{tpgb} GB activation
  exchange, and how often is it paid?
\end{qs}
\end{rightcol}
\end{document}
